\documentclass[11pt,a4paper]{article}

\usepackage[margin=1in]{geometry}
\usepackage{amsmath,amssymb}
\usepackage{graphicx}
\usepackage{booktabs}
\usepackage{csquotes}
\usepackage{hyperref}
\usepackage{natbib}
\usepackage{enumitem}
\usepackage{caption}

\hypersetup{
  colorlinks=true,
  linkcolor=black,
  citecolor=blue,
  urlcolor=blue
}

\title{\textbf{Report on ``Spatial Mapping with Gaussian Processes and Nonstationary Fourier Features''}}
\author{
  Siu Pan Lau \\
  Zhichao Zhou
}
\date{\today}

\begin{document}

\maketitle

\begin{abstract}
% Write your abstract here.
\end{abstract}

\section{Paper Information}

\begin{itemize}[leftmargin=*]
  \item \textbf{Selected paper:} ``Spatial mapping with Gaussian processes and nonstationary Fourier features''
  \item \textbf{Authors:} Jean-Francois Ton, Seth Flaxman, Dino Sejdinovic, and Samir Bhatt
  \item \textbf{Journal:} Spatial Statistics, Volume 28, pages 59--78, 2018
  \item \textbf{DOI:} \url{https://doi.org/10.1016/j.spasta.2018.02.002}
\end{itemize}

\section{Introduction}

% Introduce the selected paper/topic here.
% Use citations when referring to claims from the paper, for example: \citet{ton2018spatial}.

\section{Motivation}

% Explain why this topic is important.

\begin{quote}
``The use of covariance kernels is ubiquitous in the field of spatial statistics'' \citep[p.~59]{ton2018spatial}.
\end{quote}

\section{Background}

% Add necessary background here.

\subsection{Gaussian Processes}

% Explain Gaussian processes here.

\subsection{Stationary and Nonstationary Kernels}

% Explain stationary and nonstationary kernels here.

\subsection{Fourier Features}

% Explain Fourier features here.

\section{Method}

% Summarize the method proposed in the paper.

\section{Experiments}

% Summarize the experiments from the paper.

\begin{table}[h]
  \centering
  \caption{Summary of experiments. Fill in details from the paper.}
  \begin{tabular}{lll}
    \toprule
    Experiment & Dataset & Result/Observation \\
    \midrule
     &  &  \\
     &  &  \\
    \bottomrule
  \end{tabular}
\end{table}

\section{Discussion}

% Discuss strengths, limitations, and your group's interpretation.

\subsection{Strengths}

% Add strengths here.

\subsection{Limitations}

% Add limitations here.

\subsection{Possible Extensions}

% Add possible extensions here.

\section{Conclusion}

% Write the conclusion here.

\section*{Oral Presentation Outline}

\begin{enumerate}[label=\arabic*.]
  \item Introduction and motivation
  \item Background
  \item Method
  \item Experiments and results
  \item Discussion and conclusion
\end{enumerate}

\bibliographystyle{plainnat}
\bibliography{references}

\end{document}
